\section{Other Domains}\label{sec-other}
While substantial GraphRAG research has been investigated in previous domains, GraphRAG remains significantly underexplored in other domains such as infrastructure, biology, and scene. Therefore, we combine them into a single comprehensive section, focusing only on key studies in these fields.



\subsection{Infrastructure Graph \texorpdfstring{\emojiinfratitle}{Infrastructure Graph}}
\label{sec-infra}

%\hy{Strada-LLM: Graph LLM for traffic prediction}
%\yu{GenAI-powered Multi-Agent Paradigm for Smart Urban Mobility: Opportunities and Challenges for Integrating Large Language Models (LLMs) and Retrieval-Augmented Generation (RAG) with Intelligent Transportation Systems}

Infrastructure graphs, defined by Points of Presence (PoPs) interconnected through physical links, play a vital role in serving our daily activities across sectors such as power, water, gas, transportation, and communication. Due to the scarcity of GraphRAG research in infrastructure graphs, our review mainly focuses on graph construction methods and tasks alongside only a brief overview of recent (Graph)RAGs in this domain.

\begin{itemize}[leftmargin=*]
    \item \textbf{Power Networks}~\cite{zimmerman2010matpower, owerko2020optimal, liao2021review}: In power networks, nodes represent critical entities such as power plants, substations, transformers, and load points, while edges correspond to transmission and distribution lines facilitating the flow of electrical power across the network. Each node possesses features like power generation capacity, voltage level, load demand, geographical location, reliability metrics, and operational status, capturing the network's operational and spatial aspects. Edges are characterized by features such as transmission capacity, impedance, length, voltage level, and real-time power flow, all essential for analyzing line performance and flow efficiency. This graph-based representation enables a range of essential tasks, including transformer fault diagnosis, power outage prediction, power flow approximation, and power generation optimization.

    \item \textbf{Water Networks}~\cite{xing2022graph}: In water networks, nodes are junctions (pipe connections or users) and sources (reservoirs and tanks), and edges are pipes with water flow. Basic tasks include estimating node water heads and flow in all pipes given the water network layout, pipe characteristics, nodal demands, reservoir levels, and head measurements at limited locations~\cite{candelieri2014graph}.

    \item \textbf{Gas Networks}~\cite{wang2024multiscale, yang2023graph}: nodes represent connection points, gas sources, storage facilities, compressor stations and users, and edges edge respected natural gas pipelines, carrying user loads, diversion, and input. 

    \item \textbf{Transportation Networks}~\cite{wang2020traffic, jiang2022graph}: Road networks, where intersections serve as nodes and road connections as edges, are commonly used to model spatial dependencies for traffic flow and speed forecasting. Similarly, metro and bus networks use stations as nodes with routes as edges, capturing station topology. Region-level problems involve dividing cities into regular or irregular regions, represented as graph nodes, where spatial dependencies reflect land use patterns; regular regions often use grid partitions, while irregular ones use diverse methods like road- or zip-code-based partitions. At the road level, sensor, segment, intersection, and lane graphs capture different road elements, with nodes representing sensors, segments, intersections, or lanes. For station-level networks, nodes represent various transportation hubs (subway, bus, bike, or car-sharing stations) linked by natural connections like subway or road networks, creating an interconnected spatial dependency graph.

    \item \textbf{Communication Networks}~\cite{ferriol2023routenet, rusek2020routenet}: A computer network system typically comprises two interconnected graph layers: the logical and the physical. The logical layer represents the flow of data, where traffic is routed through multiple intermediate routers before reaching its destination. In this layer, nodes correspond to routers, and edges represent the logical paths that data packets traverse. Conversely, the physical topology layer refers to the underlying infrastructure, with Points of Presence (PoPs) from providers such as Comcast and Verizon serving as nodes and the physical fiber links between them forming the edges. Studies focused on the logical layer often model networks as graphs, where nodes represent devices (e.g., switches, routers) and edges denote links (e.g., fiber-optic cables) or traffic paths~\cite{ji2020traffic}. This graph-based representation captures both the forwarding behavior of the network (i.e., traffic interactions) and its connectivity (i.e., topological behavior), offering valuable insights for network management. For instance, optical topology-aware traffic engineering has proven effective in mitigating issues like fiber cuts and flash crowds. For the physical layer, much previous research has focused on inferring the complete as well as aligning the physical and logical layers.
\end{itemize}

To summarize, the tasks operated on infrastructure networks could be broadly encompassed as utility prediction (e.g., flow and node performance), flow simulation and generation, vulnerability analysis, and network maintenance and operation. Managing and understanding the complex physical relationships within these graphs is essential for improving infrastructure management in key areas like service delivery, function forecasting, and network optimization. In computer networks, for instance, the state of a logical traffic routing path is directly influenced by the underlying physical fiber links. Predicting critical metrics like traffic delay and jitter depends on assessing the status of all involved fiber links~\cite{ferriol2023routenet, rusek2020routenet}, as these metrics are highly interdependent. The intricate physical connections in infrastructure graphs create valuable opportunities for employing GraphRAG techniques to enhance infrastructure management and optimization. Although no existing works have explored GraphRAG in infrastructure, few works have leveraged RAG, which we briefly review in the next.

\citet{hussien2024rag} explore the use of RAG to empower automated vehicles with the ability to anticipate pedestrian and driver behaviors, including pedestrian road-crossing actions and driver lane changes. RAG retrieves explanatory documents from the JAAD and PSI datasets to provide insights into pedestrian behavior. Additionally, \citet{qian2024netbench, wang2024lens, liu2024large, bariah2024large} build on recent advances in generative models for vision and language, proposing to harness the capabilities of LLMs within computer networking. A systematic approach has been proposed to develop a foundation model for traffic tasks, such as traffic classification and generation, treating packet hexadecimal bytes as tokens. Furthermore, the potential of generative AI in advancing telecom and networking is reviewed. Specifically, \citet{kotaru2023adapting} introduces an "operator’s copilot," a natural language interface that leverages LLMs for efficient data retrieval. This interface helps manage thousands of counters and metrics, minimizing the need for specialists and accelerating issue resolution through data intelligence. These pioneering works in leveraging generative AI for computer networking lay a good foundation for rebuilding (Graph)RAG for infrastructure networks.





\subsection{Single-cell Graph \texorpdfstring{\emojibiotitle}{Biological Graph}}

Single-cell sequencing allows for the detailed analysis of molecular traits at the level of individual cells. For example, single-cell RNA sequencing (scRNA-seq)~\cite{kolodziejczyk2015technology} quantifies RNA transcript levels, offering valuable information about cell identity, developmental stages, and functional characteristics. 
Single-cell assay for transposase-accessible chromatin sequencing (scATAC-seq) ~\cite{buenrostro2013transposition} records the number of reads per accessible chromatin region, resulting in highly dimensional data matrices containing hundreds of thousands of genomic regions. 
With the explosion of the number of single-cell data, a variety of deep learning approaches ~\cite{molho2024deep, ding2024dance, tang2023single, wang2023mem, wen2023single} have been developed in recent years to tackle single-cell analysis, especially GNN methods ~\cite{lu2024graph, ding2024spatialctd, wen2022graph, shao2021scdeepsort, wangSinglecellClassificationUsing2021} have been applied to various downstream tasks. In the task of cell type annotation,
sigGCN ~\cite{wangSinglecellClassificationUsing2021} builds a gene-wise weighted adjacency matrix using data from the STRING database \cite{szklarczykSTRINGV11Protein2019} to establish a gene interaction network, with the node features representing corresponding gene expression levels. This graph is fed into a GCN-based autoencoder, which includes a convolutional layer and a max-pooling layer, followed by a flattened layer and a fully connected (FC) layer, to perform cell type annotation. 
scDeepSort \cite{shao2021scdeepsort} constructs a weighted bipartite graph where both cells and genes serve as nodes, with the edge weights representing the gene expression values for each cell-gene pair. Gene node features are derived through principal component analysis (PCA), while cell node features are obtained by aggregating the weighted features of the connected gene nodes.
The most common way to construct a single-cell graph is to build a KNN graph from single-cell data~\cite{wang2021scgnn}. To be specific, the data is first normalized to make it comparable across different cells. Next, dimensionality reduction methods, such as principal component analysis (PCA)~\cite{mackiewicz1993principal}, are employed to reduce the dataset’s high dimensionality while retaining key information. In this lower-dimensional space, distances between cells are calculated, often using Euclidean distance. Each cell’s K nearest neighbors are then determined, and a graph is formed where cells are represented as nodes, and edges connect each cell to its $K$ closest neighbors.

\textbf{Multi-omics single-cell graphs}:
% Multiome ATAC + Gene Expression; CITE-seq
Multi-omics single-cell technologies ~\cite{baysoy2023technological} integrate multiple layers of biological information—such as gene expression (RNA), chromatin accessibility (ATAC), and protein data—at the resolution of individual cells. Multiome~\cite{stuart2019integrative} refers to the simultaneous measurement of multiple molecular modalities, such as gene expression (RNA) and chromatin accessibility (ATAC), from the same single cell while CITE-seq (Cellular Indexing of Transcriptomes and Epitopes by Sequencing) ~\cite{mimitou2019multiplexed} is a multi-omics technology that allows for the simultaneous measurement of gene expression and protein surface markers at the single-cell level. 
Taking multiome dataset as an example, scMoGNN~\cite{wen2022graph} builds a heterogeneous graph consisting of cell nodes, gene nodes, and peak nodes. The edge between a cell node and a gene node indicates the expression level of that gene in the cell, while the edge between a cell node and a peak node reflects the expression level of that peak in the cell. Additionally, a threshold can be applied to filter and select the connecting edges. 
Additionally, pathways and gene activity can be incorporated as prior knowledge to establish connections between gene nodes, and connections between peaks and gene nodes respectively.
This graph representation captures both intra-modality (within each modality, like gene expression or chromatin accessibility) and cross-modality (between gene expression and chromatin accessibility) relationships, providing a comprehensive view of regulatory dynamics at the single-cell level

\textbf{Spatial transcriptomics single-cell graphs}: 
Unlike traditional transcriptomics, spatial transcriptomics retains the precise location of spot representing a cell or small group of cells within the tissue. 
Spatial transcriptomics primarily consists of three key data sources: gene expression, which captures transcript levels in specific spots; spot location, providing the spatial coordinates of each spot within the tissue; and spot images, offering visual context and morphology of the tissue at each spot.
In the task of cell type deconvolution, when building the KNN graph, GNNDeconvolver~\cite{ding2024spatialctd} integrates both cell position information and gene expression data to compute the similarity between cells. SpaGCN~\cite{hu2021spagcn} additionally incorporate spot morphology to calculate similarity between spots.
Together, constructing KNN graph with such three data sources provides a richer, more holistic view of the tissue, enabling deeper insights into cellular function, interactions, and spatial organization.


\subsection{Scene Graph \texorpdfstring{\emojiscenetitle}{Scene Graph}}


Scene graphs are a data structure designed to capture spatial and semantic relationships between objects within a scene. They are widely used in computer vision, graphics, and robotics to organize scenes, particularly when multiple objects interact, or the scene contains complex interactions. For example, SceneGraphs~\cite{he2024g} is a dataset for visual question answering, containing 100,000 scene graphs that describe the objects, attributes, and relationships within images. It presents tasks that test spatial reasoning and multi-step inference by asking users to answer open-ended questions based on textual descriptions derived from scene graphs. 

G-Retriever~\cite{he2024g} processes Scene Graphs by first parsing JSON data to index objects and attributes, allowing for efficient retrieval of relevant nodes and edges based on the query. It then constructs a subgraph with the necessary scene information, filtering out unrelated details. Finally, G-Retriever uses this subgraph along with an LLM to generate answers, leveraging spatial and semantic relationships for accurate reasoning and response generation.

P-RAG~\cite{xu2024p} engages with the environment through agents, building and updating a database of historical trajectories that guide the agent’s actions. Initially, the task’s goal instruction is provided. Before each action, the agent captures an observation image reflecting the current state, which is then converted into a scene graph format to facilitate processing by LLMs. By creating and retrieving scene graph context, P-RAG enhances the LLM's ability to accurately interpret and navigate complex scenes, making this approach especially useful in planning embodied everyday tasks.



\subsection{Random Graph \texorpdfstring{\emojirandomtitle}{Random Graph}}


Random graphs are a foundational concept in network theory~\citep{newman2018networks, janson2011random} and are widely used for modeling and studying complex networks in fields such as computer science, physics, biology, and social sciences. They are constructed based on probabilistic rules, leading to various possible structures that reflect the diversity seen in real-world networks. These random graphs can be used to analyze the retrieval, organization, and generation processes within GraphRAG.

There are many models to generate random graphs, such as Erdős-Rényi Model (ER Model)~\citep{erdds1959random}, Watts–Strogatz model~\citep{watts1998collective}, Barabási–Albert model~\cite{albert2002statistical}, Random geometric graph~\cite{cannings2005random}, scale-free networks~\citep{barabasi1999emergence} and stochastic block model~\cite{holland1983stochastic}. Additionally, various graph structures, such as Path Graphs, Complete Graphs, Star Graphs, and Barbell Graphs, can be generated to meet specific analytical needs.

Random graph generation models, such as Contextual stochastic block models (CSBM)~\citep{deshpande2018contextual} has been widely leveraged in GNNs analysis~\citep{baranwal2021graph, ma2021homophily, mao2024demystifying, han2024node}. Additionally, random graphs are increasingly applied to examine the behavior of LLMs. For instance, \citet{fatemi2023talk} leverage Erdős-Rényi graphs, scale-free networks, Barabási–Albert, stochastic block model, star, path and complete graph generators to test LLM performance on various graph reasoning tasks, such as edge existence, node degree, node/edge count, connected nodes and cycle check.
GraphLLM~\citep{chai2023graphllm} generates random graphs in various input formats to assess LLM performance on graph reasoning tasks, including substructure counting, maximum triplet sum, shortest path calculation, and bipartite graph matching. In this approach, a graph transformer encodes the graph, and the embedding-fusion method is used for generation. \citet{bachmann2024pitfalls} use star graphs to investigate limitations within the next-token prediction paradigm. Other studies, such as \citet{dai2024revisiting}, \citet{wang2024can}, \citet{guo2023gpt4graph}, and \citet{luo2024graphinstruct}, explore the graph reasoning abilities of LLMs across various tasks.


